\documentclass[a4paper, serif, 10pt]{moderncv}

%% moderncv
% color and style
\moderncvcolor{black}
\moderncvstyle{banking}

%% page layout/margins
\usepackage[top=2.5cm, bottom=2.5cm, left=1.5cm, right=1.5cm]{geometry}

\nopagenumbers{}

%% Babel config for Indonesian (Bahasa Indonesia) language
% ref:
% - https://latex3.github.io/babel/guides/locale-indonesian.html
\usepackage[indonesian]{babel}

%% fontspec
\usepackage{fontspec}

\IfFontExistsTF{Linux Libertine}{
  \setmainfont[Ligatures={TeX, Common}, Numbers=OldStyle]{Linux Libertine}
}{}
\IfFontExistsTF{Linux Libertine O}{
  \setmainfont[Ligatures={TeX, Common}, Numbers=OldStyle]{Linux Libertine O}
}{}

%% CTeX
% CJK support, used to show CJK characters in the resume
%
% - heading=false: disable ctex heading
% - punct=kaiming: use kaiming punctuations styles for CJK
% - scheme=plain: use plain scheme, do not override \normalsize font size
% - space=auto: space settings for CJK characters
%
% CJK fonts are resolved per locale language, e.g. Noto Serif/Sans CJK JP
% for Japanese resumes, with fallback to the next region if not installed
%
% ref:
% - http://ctan.mirrorcatalogs.com/language/chinese/ctex/ctex.pdf
\usepackage[UTF8, heading=false, punct=kaiming, scheme=plain, space=auto]{ctex}

\IfFontExistsTF{Noto Serif CJK SC}{
  \setCJKmainfont{Noto Serif CJK SC}
}{}
\IfFontExistsTF{Noto Sans CJK SC}{
  \setCJKsansfont{Noto Sans CJK SC}
}{}

%% line spacing
\usepackage{setspace}
\setstretch{1.125}

%% URL styling - use normal text instead of monospace
\urlstyle{same}

%% Patch \httplink and \httpslink to handle full URLs
% The original moderncv macros always prepend http:// or https:// to URLs.
% This causes issues when the URL already has a protocol (e.g., https://example.com)
% resulting in malformed URLs like https://https://example.com.
% This patch checks if the URL already contains "://" and uses it directly if so.
% Note: We use \str_set:Nx (x-expansion) to expand macros like \@homepage first.
\ExplSyntaxOn
\str_new:N \g_yamlresume_url_str

\RenewDocumentCommand{\httpslink}{O{}m}{
  \str_set:Nx \g_yamlresume_url_str {#2}
  \str_if_in:NnTF \g_yamlresume_url_str {://}
    { \url{#2} }
    { \url{https://#2} }
}

\RenewDocumentCommand{\httplink}{O{}m}{
  \str_set:Nx \g_yamlresume_url_str {#2}
  \str_if_in:NnTF \g_yamlresume_url_str {://}
    { \url{#2} }
    { \url{http://#2} }
}
\ExplSyntaxOff

\name{Rizky Pratama}{}
\title{Manajer Operasional Berpengalaman di Ritel dan Logistik}
\phone[mobile]{+62 812 3456 7890}
\email{rizky.pratama@example.co.id}
\homepage{https://www.rizkypratama.id}

\address{Jl. Jenderal Sudirman No. 45, Jakarta Selatan, DKI Jakarta, Indonesia, 12190}{}{}

\extrainfo{{\small \faLinkedin}\ \href{https://www.linkedin.com/in/rizkypratama}{@rizkypratama} {} {} {} • {} {} {} 
{\small \faTwitter}\ \href{https://twitter.com/rizkypratama}{@rizkypratama}}

\begin{document}

\maketitle

\section{Informasi Dasar}

\cvline{}{Manajer operasional dengan pengalaman lebih dari 8 tahun dalam mengelola proses ritel, distribusi, dan logistik. Terbukti mampu meningkatkan efisiensi operasional, menurunkan biaya, dan memimpin tim lintas fungsi untuk mencapai target KPI.
%
\begin{itemize}
\item Ahli dalam perencanaan kapasitas, pengendalian persediaan, dan perbaikan proses berkelanjutan.
\item Berpengalaman memimpin tim hingga 60 orang di beberapa lokasi operasional.
\item Kuat dalam analisis data operasional, penyusunan anggaran, dan manajemen pemasok.
\end{itemize}}
\section{Pendidikan}

\cventry{Sep 2011–Jul 2015}
        {Sarjana, Manajemen Operasional, Nilai: 3.65}
        {Universitas Indonesia}
        {\url{https://www.ui.ac.id}}
        {}
        {\begin{itemize}
\item Menempuh studi dengan fokus pada efisiensi proses, pengendalian persediaan, dan manajemen rantai pasok.
\item Menyelesaikan proyek akhir tentang optimasi rute distribusi untuk menurunkan biaya logistik.
\item Aktif dalam organisasi kemahasiswaan sebagai koordinator acara operasional.
\end{itemize}
\textbf{Mata Kuliah}: Manajemen Rantai Pasok, Perencanaan dan Pengendalian Produksi, Analisis Proses Bisnis, Manajemen Kualitas, Manajemen Proyek, Riset Operasi}
\section{Pengalaman Kerja}

\cventry{Mar 2021–Sekarang}
        {Manajer Operasional Regional}
        {PT Nusantara Retail Mandiri}
        {\url{https://www.nusantararetail.co.id}}
        {}
        {\begin{itemize}
\item Memimpin operasional 18 gerai dan 3 gudang regional dengan total 60 staf.
\item Menyusun rencana operasional bulanan dan tahunan, termasuk anggaran, target KPI, dan alokasi sumber daya.
\item Meningkatkan produktivitas tim sebesar 22\% melalui standarisasi proses dan program pelatihan berkelanjutan.
\item Menurunkan biaya logistik regional sebesar 15\% melalui negosiasi vendor dan optimasi rute pengiriman.
\item Berkolaborasi dengan tim penjualan, keuangan, dan rantai pasok untuk memastikan ketersediaan produk dan kepatuhan SOP.
\end{itemize}
\textbf{Kata Kunci}: Perencanaan Operasional, Efisiensi Proses, Manajemen Tim, KPI, Logistik}

\cventry{Jun 2017–Feb 2021}
        {Asisten Manajer Operasional}
        {PT Sentra Logistik Cepat}
        {\url{https://www.sentralogistik.co.id}}
        {}
        {\begin{itemize}
\item Mengelola operasional harian pergudangan dan distribusi untuk wilayah Jabodetabek.
\item Mengawasi tim 25 orang dan memastikan pencapaian target pengiriman tepat waktu di atas 98\%.
\item Mengembangkan dashboard KPI untuk memantau produktivitas, biaya per pengiriman, dan tingkat kerusakan barang.
\item Memimpin inisiatif perbaikan proses yang mengurangi waktu bongkar muat sebesar 30\%.
\end{itemize}
\textbf{Kata Kunci}: Pergudangan, Distribusi, Dashboard KPI, Perbaikan Proses, Manajemen Tim}

\cventry{Agu 2015–Mei 2017}
        {Supervisor Operasional}
        {PT Karya Distribusi Utama}
        {\url{https://www.karyadistribusi.co.id}}
        {}
        {\begin{itemize}
\item Mengawasi kegiatan inventaris, penerimaan barang, dan pengiriman ke pelanggan ritel.
\item Menyusun jadwal kerja harian untuk 15 staf operasional dan memastikan kepatuhan terhadap prosedur keselamatan.
\item Melakukan analisis ABC untuk mengoptimalkan persediaan dan mengurangi stok mati sebesar 18\%.
\end{itemize}
\textbf{Kata Kunci}: Pengendalian Persediaan, Jadwal Kerja, Keselamatan Kerja, Analisis ABC, Ritel}
\section{Bahasa}

\cvline{Indonesia}{Bahasa Ibu / Bilingual \hfill \textbf{Kata Kunci}: Bahasa ibu}
\cvline{Inggris}{Tingkat Profesional Penuh \hfill \textbf{Kata Kunci}: TOEFL 580, Presentasi Bisnis}
\section{Keahlian}

\cvline{Manajemen Operasional}{Ahli \hfill \textbf{Kata Kunci}: Perencanaan Kapasitas, Pengendalian Persediaan, Analisis Proses, Manajemen Kualitas}
\cvline{Kepemimpinan}{Mahir \hfill \textbf{Kata Kunci}: Manajemen Tim, Coaching, Delegasi, Komunikasi Lintas Fungsi}
\cvline{Analisis Data}{Mahir \hfill \textbf{Kata Kunci}: Excel, Power BI, SQL, Dashboard KPI}
\cvline{Manajemen Rantai Pasok}{Mahir \hfill \textbf{Kata Kunci}: Logistik, Pengadaan, Distribusi, Pergudangan}
\cvline{Manajemen Proyek}{Menengah \hfill \textbf{Kata Kunci}: Agile, Scrum, Perencanaan Proyek, Manajemen Risiko}
\section{Penghargaan}

\cventry{Des 2022}
        {PT Nusantara Retail Mandiri}
        {Penghargaan Kinerja Terbaik}
        {}
        {}
        {Diberikan kepada manajer operasional dengan pencapaian KPI tertinggi dan kontribusi signifikan terhadap efisiensi biaya regional.}

\cventry{Agu 2019}
        {PT Sentra Logistik Cepat}
        {Inisiatif Efisiensi Biaya}
        {}
        {}
        {Mengakui keberhasilan program pengurangan biaya distribusi melalui optimasi rute dan negosiasi vendor.}
\section{Sertifikat}

\cventry{Mar 2020}
        {Asosiasi Logistik Indonesia}
        {Sertifikasi Manajemen Rantai Pasok Profesional}
        {\url{https://www.asosiasilogistik.id/sertifikasi}}
        {}
        {}

\cventry{Nov 2021}
        {Lembaga Produktivitas Nasional}
        {Lean Six Sigma Green Belt}
        {\url{https://www.lpn.go.id/lean-six-sigma}}
        {}
        {}

\cventry{Feb 2018}
        {Kementerian Ketenagakerjaan RI}
        {Ahli K3 Umum}
        {\url{https://www.kemnaker.go.id}}
        {}
        {}
\section{Publikasi}

\cventry{Apr 2023}
        {Optimalisasi Persediaan dengan Analisis ABC dan Forecast Permintaan}
        {Majalah Logistik Indonesia}
        {\url{https://www.majalahlogistik.id/artikel/optimalisasi-persediaan}}
        {}
        {\begin{itemize}
\item Membahas pendekatan praktis untuk mengklasifikasikan persediaan berdasarkan nilai dan tingkat perputaran.
\item Menjelaskan integrasi forecast permintaan dengan analisis ABC untuk mengurangi biaya penyimpanan.
\item Studi kasus pada jaringan ritel regional dengan lebih dari 500 SKU.
\end{itemize}}
\section{Referensi}

\cventry{\emaillink[dewi.lestari@nusantararetail.co.id]{dewi.lestari@nusantararetail.co.id}}
        {Atasan Langsung di PT Nusantara Retail Mandiri}
        {Dewi Lestari}
        {+62 811 2233 4455}
        {}
        {Rizky adalah manajer operasional yang berorientasi pada hasil dan mampu memimpin tim besar dengan jelas. Ia konsisten dalam mencapai target dan membangun proses yang lebih efisien.}

\cventry{\emaillink[budi.santoso@sentralogistik.co.id]{budi.santoso@sentralogistik.co.id}}
        {Mantan Atasan di PT Sentra Logistik Cepat}
        {Budi Santoso}
        {+62 813 9988 7766}
        {}
        {Selama bekerja bersama, Rizky menunjukkan kemampuan analitis yang kuat dan komitmen tinggi terhadap perbaikan berkelanjutan. Ia layak dipercaya untuk peran kepemimpinan operasional.}
\section{Proyek}

\cventry{Jan 2022–Nov 2023}
        {Implementasi sistem WMS dan otomasi pelaporan KPI di 12 gudang regional.}
        {Transformasi Digital Operasional Gudang}
        {\url{https://www.nusantararetail.co.id/proyek/digital-gudang}}
        {}
        {\begin{itemize}
\item Memimpin implementasi WMS untuk meningkatkan akurasi stok dan visibilitas persediaan.
\item Mengotomasi pelaporan KPI harian sehingga waktu pelaporan berkurang dari 4 jam menjadi 30 menit.
\item Melatih 80 staf gudang dan supervisor dalam penggunaan sistem baru.
\end{itemize}
\textbf{Kata Kunci}: WMS, Otomasi, KPI, Manajemen Perubahan}

\cventry{Mar 2019–Des 2020}
        {Inisiatif lintas fungsi untuk menurunkan biaya distribusi tanpa mengurangi kualitas layanan.}
        {Program Pengurangan Biaya Logistik}
        {\url{https://www.sentralogistik.co.id/program/efisiensi-biaya}}
        {}
        {\begin{itemize}
\item Menganalisis data pengiriman dan mengidentifikasi peluang konsolidasi rute.
\item Bernegosiasi dengan vendor transportasi untuk mendapatkan tarif yang lebih kompetitif.
\item Mencapai penghematan biaya logistik sebesar 12\% dalam 18 bulan.
\end{itemize}
\textbf{Kata Kunci}: Efisiensi Biaya, Rute Distribusi, Negosiasi, Analisis Data}
\section{Minat}

\cvline{Olahraga}{Lari, Bersepeda, Bulu Tangkis}
\cvline{Pengembangan Diri}{Membaca, Public Speaking, Mentoring}
\section{Kegiatan Sukarela}

\cventry{Jan 2020–Des 2023}
        {Koordinator Relawan Distribusi Bantuan}
        {Yayasan Peduli Logistik Indonesia}
        {\url{https://www.pedulilogistik.org}}
        {}
        {\begin{itemize}
\item Mengoordinasikan distribusi bantuan logistik ke daerah terdampak bencana di Jawa Barat.
\item Mengelola 30 relawan dan berkoordinasi dengan pemerintah daerah serta mitra donor.
\item Menyusun prosedur penerimaan, penyimpanan, dan penyaluran barang agar transparan dan akuntabel.
\end{itemize}}
    
\end{document}